% ---
% title: Infinite Dimension ™
% date: 2026-04-07
% tag: Math
% excerpt: 变分法的一页笔记：泛函、函数空间、第一变分，以及为什么无穷维里的极值问题需要另一套语言。
% ---
\href{https://www.cnblogs.com/qizhou/p/18564508}{泛函微分，是为变分。}

\section{Calculus of Variations}

Suppose \(J\) is a functional on the function space \(\mathcal{F}[a,b]\), with the form

\[
J[y] = \int_{a}^{b}L(x,y(x),y'(x))dx \tag{1}
\]

where \(y\) is an admissible function with a continuous first derivative, subject to \(y(a)=A, \quad y(b)=b\).

The classical problem is to find the \(y\) that minimizes \(J\).

\hrule

A functional is easier to understand next to the finite-dimensional case.
In finite dimensions, a function \(f:\mathbf{R}^n \to \mathbf{R}\) changes at a point \(x_0\) along a direction \(v\) according to the directional derivative, i.e.,
\(\nabla f(x_0)\cdot v\).

The analogous concept for a functional is the variation. We denote by \(\delta y(x)\) an admissible perturbation of the function \(y(x)\). To satisfy the endpoint constraints, we require \(\delta y(a)=\delta y(b)=0\).

Then the first-order variation of the functional \(J\) at the function \(y\) along the direction \(\delta y\) is defined by

\[
\delta J[y;\delta y]:=\left.\frac{d}{d\epsilon}J[y+\epsilon\delta y]\right|_{\epsilon=0}. \tag{2}
\]

Assume the first variation is linear and continuous in \(\delta y\).
Then there exists a quantity \(\frac{\partial J}{\partial y(x)}\), called the functional derivative of \(J\), such that

\[
\delta J[y;\delta y]=\int_a^b \frac{\partial J}{\partial y(x)}\,\delta y(x)\,dx. \tag{3}
\]

For the functional

\[
J[y]=\int_a^b L(x,y,y')\,dx,
\]

the first variation \(\delta J[y;\delta y]\) is the differential of \(J\) in the direction \(\delta y\). It can be computed explicitly as

\begin{align}
\delta J[y;\delta y]
&=\left.\frac{d}{d\epsilon}J[y+\epsilon\delta y]\right|_{\epsilon=0} \\
&=\left.\frac{d}{d\epsilon}\int_a^b
L\big(x,\,y+\epsilon\delta y,\,(y+\epsilon\delta y)'\big)\,dx\right|_{\epsilon=0} \\
&=\left.\int_a^b
\frac{d}{d\epsilon}
L\big(x,\,y+\epsilon\delta y,\,y'+\epsilon\delta y'\big)\,dx\right|_{\epsilon=0} \\
&=\int_a^b
\left(
\frac{\partial L}{\partial y}\,\delta y
+
\frac{\partial L}{\partial y'}\,\delta y'
\right)dx.
\end{align} \tag{4}

The term \(\int_a^b \frac{\partial L}{\partial y'}\delta y'\,dx\) integrates by parts:

\begin{align}
\int_a^b \frac{\partial L}{\partial y'}\delta y'dx &=\int_a^b \frac{\partial L}{\partial y'}d\delta y  \\
&=\frac{\partial L}{\partial y'}\delta y \large|_a^b - \int_a^b \frac{d }{dx} \frac{\partial L}{\partial y'} \delta ydx \\
&= -\int_a^b \frac{d }{dx} \frac{\partial L}{\partial y'} \delta y dx 
\end{align} \tag{5}

Substituting (5) into (4) gives the functional derivative:

\[
\begin{aligned}
\delta J[y;\delta y]
&= \int_a^b \left(\frac{\partial L}{\partial y}\delta y+\frac{\partial L}{\partial y'}\delta y'\right)dx \\
&= \int_a^b \frac{\partial L}{\partial y}\delta y\,dx
+ \int_a^b \frac{\partial L}{\partial y'}\frac{d}{dx}(\delta y)\,dx \\
&= \int_a^b \frac{\partial L}{\partial y}\delta y\,dx
+ \left[\frac{\partial L}{\partial y'}\delta y\right]_a^b
- \int_a^b \frac{d}{dx}\left(\frac{\partial L}{\partial y'}\right)\delta y\,dx \\
&= \left[\frac{\partial L}{\partial y'}\delta y\right]_a^b
+ \int_a^b \left(\frac{\partial L}{\partial y}-\frac{d}{dx}\frac{\partial L}{\partial y'}\right)\delta y\,dx.
\end{aligned}
\]

Thus,

\[
\frac{\partial J}{\partial y(x)} = \int_a^b\big ( \frac{\partial L}{\partial y} - \frac{d }{dx} \frac{\partial L}{\partial y'} \big ) dx
\]

A necessary condition for \(y\) to extremize \(J\) is \(\frac{\partial J}{\partial y(x)}=0\quad \forall x \in [a,b]\).
